\section{Introduction}
\label{sec:intro}

Language models are increasingly deployed in specialized domains---legal analysis, biomedical research, scientific literature---where the target distribution differs substantially from pretraining data. The standard approach to bridging this gap is domain-specific fine-tuning, which requires additional data, compute, and engineering effort for each new domain. A more appealing alternative would be to pretrain models that adapt \emph{rapidly} to new distributions, reducing the cost of domain specialization.

We take inspiration from an unlikely source: how distribution shifts during training might serve as a feature rather than a bug. Meta-learning methods train models across diverse tasks precisely so they can adapt quickly to new ones~\cite{finn2017maml}. Curriculum learning research has shown that training data ordering affects what models learn and how efficiently they learn it~\cite{bengio2009curriculum, fan2023irreducible, chen2023skillit}. Yet no prior work has tested whether deliberately introducing domain-level distribution shifts during pretraining can improve a model's ability to adapt to new domains.

{\bf Our main question is:} does training a language model with phased domain introduction---withholding certain domains initially, then introducing them---produce a model that adapts faster to unseen distributions compared to standard uniform pretraining?

We propose a \emph{distribution shift curriculum} (\dsc), a simple two-phase pretraining strategy. In Phase~1, the model trains on a subset of ``core'' domains. In Phase~2, previously withheld ``shift'' domains are introduced alongside the core domains. This controlled distribution shift forces the model to integrate new data distributions into its existing representations---a form of implicit meta-learning at the domain level.

We test \dsc on a 30M-parameter GPT-2 model using 13 domains from \thepile~\cite{gao2020pile}. Our experiments reveal three findings. First, the curriculum model adapts 2.8$\times$ faster to shift domains during a post-training adaptation phase (9.1 vs.\ 3.3 perplexity improvement over 500 steps). Second, \dsc unexpectedly improves performance on core domains: the curriculum model outperforms the uniform baseline on 7 of 9 core domains, with 5 reaching statistical significance ($p < 0.05$). Third, the curriculum model exhibits 26\% less catastrophic forgetting on held-out domains during adaptation. However, \dsc does not improve zero-shot transfer to completely unseen domains, and models end training with higher perplexity on shift domains they saw for fewer steps.

Our contributions are:
\begin{itemize}[leftmargin=*,itemsep=0pt,topsep=0pt]
    \item We propose \dsc, a domain-withholding curriculum for language model pretraining that creates controlled distribution shifts to improve adaptation speed.
    \item We conduct controlled experiments (3 seeds, 13 domains) showing that \dsc produces models that adapt 2.8$\times$ faster to new domains while improving core-domain performance by 5.3\%.
    \item We identify an unexpected benefit of domain concentration: training on fewer domains in Phase~1 yields stronger representations that generalize better within the core domain set.
    \item We provide detailed per-domain analysis and ablations that characterize when and why distribution shift curricula help or hurt.
\end{itemize}
